\section{Einleitung}

\subsection{Motivation und Problemstellung}

In automatisierten Industrieprozessen ist eine leistungsfähige und strukturierte Infrastruktur erforderlich, um einen reibungslosen und effizienten Ablauf zu gewährleisten. Solche Systeme sind hochgradig diversifiziert, da unterschiedliche Teilprozesse und Aufgaben gemeinsam zur Realisierung eines übergeordneten Produktions- oder Logistikablaufs beitragen.

Fehlen geeignete Strukturen oder sind diese unzureichend ausgelegt, kann dies zu Produktionsfehlern, Lieferverzögerungen oder ungeplanten Störungen innerhalb des Prozesszyklus führen. Eine zentrale Rolle nimmt in diesem Zusammenhang die \textbf{Logistik} ein. Logistische Prozesse tragen maßgeblich zur Stabilität, Kontinuität und Flexibilität industrieller Abläufe bei.

\subsection{Bedeutung der Lagerlogistik}

Innerhalb der Fertigung müssen Materialien, Bauteile und Produkte bedarfsgerecht bereitgestellt werden. Dies erfolgt über strukturierte logistische Prozesse, die eine zeit- und ortsgerechte Bereitstellung sicherstellen. Dabei spielt insbesondere die \texttt{Lagerlogistik} eine entscheidende Rolle.

In diesem Bereich werden Materialien organisiert gelagert, verwaltet und für nachgelagerte Prozesse verfügbar gemacht. Häufig erfolgt dies in Form von Hochregallagern, die eine platzsparende und automatisierte Lagerhaltung ermöglichen.

\subsection{Versuchsaufbau und bisherige Arbeiten}

Zur praxisnahen Abbildung solcher Prozesse steht an der \textbf{DHBW Ravensburg – Campus Friedrichshafen} ein didaktisches Modell eines Hochregallagers zur Verfügung. Dieses System ermöglicht es Studierenden, praktische Erfahrungen im Bereich der automatisierten Logistik zu sammeln.

In der vorangegangenen Arbeit im Modul \textbf{T3\_3100} wurden bereits Anpassungen und Erweiterungen an diesem System vorgenommen. Der Fokus lag dabei insbesondere auf sicherheitsrelevanten Aspekten sowie auf dem Prozess des \textbf{Umlagerns}.\newpage

\subsection{Zielsetzung der Arbeit}

Die vorliegende Arbeit baut auf diesen Grundlagen auf und erweitert das bestehende System um eine realitätsnahe Funktionalität:\\
Die \textbf{auftragsbasierte Steuerung des Auslagerprozesses}.\\
In modernen industriellen Lagerverwaltungssystemen stellt die auftragsbasierte Prozesssteuerung einen wesentlichen Bestandteil dar, um Effizienz, Nachvollziehbarkeit und Flexibilität zu gewährleisten.

Ziel dieser Arbeit ist es, ein entsprechendes Steuerungskonzept zu entwickeln und auf das vorhandene Modell zu übertragen. Dabei wird insbesondere die Umsetzung einer auftragsgesteuerten Auslagerung unter Einsatz von \textbf{Python} betrachtet. Die Integration der Softwarelösung ermöglicht eine flexible Erweiterung der bestehenden Steuerungsarchitektur sowie eine verbesserte Visualisierung und Verarbeitung von Prozessdaten.\\